\documentclass[../main.tex]{subfiles}
\begin{document}
%==========================================TIÊU ĐỀ TẬP========================================
\begin{table}[h]
  \begin{adjustwidth}{-1cm}{}
    \begin{tabular}{>{\centering\arraybackslash}p{6cm}|>{\raggedright\arraybackslash}p{12.5cm}}
      \multirow{5}{6cm}
      % Cột bên trái: ÉPISODE và số tập
      {\\[-40pt]
        \flaregothic\fontsize{20pt}{20pt}\selectfont \centering
        \color{eptype}ÉPISODE\color{black}\\
        \burbank\fontsize{72pt}{72pt}\selectfont
        \color{epnum}114\color{black}\\[6pt]
        \cafeteria\fontsize{20pt}{20pt}\selectfont\color{epnum}(9-5)\color{black}
        \\[18pt]
        % \flaregothic\fontsize{14pt}{14pt}\selectfont
        % \color{epattr}BIENVENUE, \uline{\color{Banpire} Mel} !
        \color{black}
      }
       &
      % Dòng thứ nhất: tên tập tiếng Nhật
      {\fotskip\fontsize{18pt}{18pt}\selectfont \ruby{私}{わたし}に\ruby{関}{か}わるな
      }                 \\[6pt]
      % Dòng thứ hai: tên tập tiếng Anh
       & {\cafeteria\fontsize{18pt}{18pt}\selectfont Stay Out of My Lane}                   \\[4pt]
      % Dòng thứ ba: tên tập tiếng Việt
       & {\vnmsans\fontsize{18pt}{18pt}\selectfont Bóng hồn Aqua}                           \\[8pt]
      % Dòng thứ tư: Nhân vật xuất hiện
       & {\sen{\fontsize{14pt}{14pt}\selectfont Track used: Sky (\textnormal{\ruby{空}{そら}})}
         }
      \\[8pt]
      % Dòng cuối cùng: Thời gian diễn ra của tập (thường sẽ là ngày phát hành)
       & {\continuum\fontsize{16pt}{16pt}\selectfont Ngày phát hành: 18/04/2021}            \\[18pt]
    \end{tabular}
  \end{adjustwidth}
\end{table}
%=========================================TRUYỆN CHÍNH========================================
\color{black}

\pov{Hikawa Kuon}

Một ngày sau sinh nhật tôi.

Không khí náo nhiệt đã lùi lại phía sau. Mọi người đã trở lại với guồng quay công việc thường nhật, như thể hôm qua chưa từng có bữa tiệc ồn ào đến vậy.

Nhưng như thường lệ ở văn phòng \hll, ``bình thường'' chỉ là một điều gì đó rất viển vông.

Ngay lúc này đây, Korone-chan đang vui vẻ\dots\ véo má Akuta, người bằng cách nào đó đã ngất xỉu và đang nằm bẹp trên sàn, mà tôi sẽ kể lại sau. Cái đuôi của em ấy thì vẫy liên tục như không có chuyện gì nghiêm trọng cả.
\img{9-5a}

``Không lẽ bà ấy hết năng lượng rồi sao?'' nàng Khuyển hỏi, hai tay vẫn không ngừng véo đối phương với sự thích thú.

Tôi thở dài, đứng khoanh tay sau lưng, ngán ngẩm nhìn cảnh tượng trước mắt: ``Em cứ làm như vừa rồi người ta không xỉu mới lạ đấy\dots''

Tôi còn chưa kịp giải thích thì *RẦM!* Một tiếng động cực mạnh vang lên sau lưng. Tôi và Korone đồng loạt quay đầu lại.

Đó là Kanata.

\gif{9-5b}{2}{74}

Cô nàng Thiên sứ đang nằm sõng soài giữa một mảnh trần nhà vừa bị rơi. Tường vỡ, bụi bốc lên, mảnh trần vỡ toạc nằm chỏng chơ.

Kanata ngồi dậy, vừa xoa đầu vừa nhăn mặt: ``Ui da\dots\ đau ghê\dots''

Rồi vừa thấy chúng tôi đang nhìn mình chằm chằm, em ấy lập tức dựng người lên, chống tay và vẫy vẫy: ``À-À, chảo hai người\~{}!''

Câu chào ngốc nghếch không cứu vãn được tình hình. Korone đứng nghiêm, không nói gì, còn tôi thì chỉ biết thở dài, khoanh tay nhìn cô ấy bằng ánh mắt bất lực.

Chúng tôi cùng đồng thanh, không hẹn mà nên:
``Rồi, xong đời luôn.''

Kana-chan ngơ ngác. Cô nhìn qua nhìn lại giữa chúng tôi, hoàn toàn không hiểu. Chỉ đến khi cúi xuống, và thấy mình đang ngồi hẳn lên mảnh trần nhà vỡ toang\dots\ em ấy mới hiểu rõ vấn đề, khẽ buột miệng: ``À\dots''

Ngay lúc ấy, cánh cửa văn phòng bật mở.

A-san xuất hiện, hai tay đang bưng một thùng tài liệu. Nhưng ngay khi ngẩng đầu lên và thấy cái trần nhà toác một lỗ, ánh mắt chị ấy lập tức đông cứng. Chiếc hộp rơi khỏi tay trong vô thức, tài liệu tung tóe ra sàn.

Cả tôi cùng hai cô gái đứng nhìn A-san đang dần mất kiểm soát. Không ai dám nói gì. Ngay khi một trong chúng tôi định lên tiếng, chị ấy liền ôm đầu, khuỵu xuống thảm thiết: ``\textbf{Không!!}''

Tôi hoảng hốt lao tới đỡ cô, sợ cô ngã thật thì hậu quả sẽ còn khủng khiếp hơn. Còn Koro-chan thì vẫn giữ nét mặt ngây thơ pha chút\dots\ thương hại.

Nàng Khuyển và tôi đồng thanh nói lớn với kẻ vừa phá hỏng trần nhà ấy: ``Tôi đã bảo rồi, đi đời luôn.'' / ``Liệu mà làm cách nào sửa cái trần đi nha.''

Kanata không dám cãi lại. Với vẻ mặt buồn rầu và sợ sệt như bị trừ điểm hạnh kiểm, em ấy miễn cưỡng đi bắc thang.

Chẳng biết em ấy đào đâu ra một ít xi măng, nhặt mấy miếng trần nhà bị vỡ ở dưới, cuộn băng dính khổ to, rồi\dots\ dùng cả sức mạnh thiên sứ để ``tạm vá'' lại cái lỗ vừa tạo ra.

\img{9-5c}

Không cần ai nhắc nhở, vì đơn giản: em ấy không muốn đền tiền. Tôi rất muốn nói rằng ``Của rẻ là của ôi đấy!'' nhưng dám chắc em ấy cũng chẳng muốn nghe đâu.

Thôi thì, em ấy dần cũng sẽ thấy thôi.

\ct

Sau một hồi loay hoay trên thang, Kana-chan cuối cùng cũng dán được mấy mảng trần tạm bợ kia trở lại chỗ cũ. Xi măng bám lổm chổm, băng dính vắt chéo lung tung như mạng nhện, trần nhà bây giờ trông chẳng khác gì một bản vá lỗi thời gian thực. Em ấy vừa trèo xuống, vừa lau mồ hôi trán, thở hắt ra như thể vừa làm xong một việc vĩ đại.

Tôi ngẩng đầu nhìn cái thành quả ``tạm ổn'' ấy và không nhịn được hỏi: ``Thế cái đó có ổn không vậy? Anh thấy trông vẫn còn\dots\ lủng lẳng lắm đấy.''

Nhưng em ấy chỉ xua tay, mỉm cười méo xệch: ``Em thấy thế cũng được rồi mà\dots\ Cứ để như vậy đã rồi tính sau\dots\''

Nói rồi, em quay đầu lại, hướng ánh mắt về phía tiền bối vẫn đang chơi đùa như chẳng có chuyện gì xảy ra với Aqua, người vẫn đang nằm bất tỉnh nhân sự. Chẳng mất nhiều thời gian để em lại túa mồ hôi hột: ``Mà tiện thể thì\dots\ Chị đang làm cái gì vậy\dots?''

Korone-chan vẫn giữ nguyên nụ cười toe toét, tay tiếp tục xoa má Aqua-chan một cách vui vẻ như đang cưng nựng chó cưng. ``Aqua-chan bất tỉnh rồi đó, em thấy không?''

\img{9-5d}

``Ủa? Không phải là đang ngủ ạ?'' Thiên sứ đổ mồ hôi hột khi thấy tiền bối của cô không hề giấu giếm gì.

Tôi khoanh tay, chán chẳng buồn nói, chỉ liếc sang người đang vác nàng Hầu gái lên đầy ngụ ý: ``Nếu không vì \emph{ai đó} làm ra nông nỗi này. Với lại, nếu đang ngủ thì đáng ra em ấy phải thở mạnh, cựa quậy hay gì đó chứ?''

``Vừa nãy em vừa phi một cái *RẦM!* mạnh thế mà bà ấy cũng chẳng phản ứng gì,'' nàng Khuyển nói tiếp, ngón tay em chạm mạnh vào má của nàng Hầu gái đen đủi ấy.

``Ngủ say đến mức bất tỉnh nhân sự luôn,'' tôi mỉa mai.

\dots

Một thoáng im lặng trôi qua, rồi em ấy lại hỏi, lần này là hỏi cả hai chúng tôi: ``Anh chị có biết tại sao lại thế này không?''

Koro-chan bắt đầu gãi đầu, giả bộ nhớ lại điều gì đó\dots\ rồi ngay lập tức lắc đầu, tỏ ra vô tội như thiên thần: ``Không biết!''

Tất nhiên là em Thiên sứ không dễ bị lừa. Em híp mắt lại, chỉ tay vào tiền bối: ``Xạo ke! Chị biết rõ là đằng khác!''

Ngay sau đó, em quay sang tôi, ra vẻ nghiêm túc như một điều tra viên hỏi cung nhân chứng: ``Anh ở đây với Korone-senpai, anh cũng phải biết gì chứ?''

Tôi thì vốn dĩ chẳng hề có ý định trốn tránh gì cả, nên cũng không tiếc mắc mà gật đầu. ``Được rồi, chuyện là thế này\dots'' tôi bắt đầu kể bằng giọng điệu chán không tả nổi.

\begin{center}
  \uline{Hồi tưởng.}
\end{center}

Khi tôi vừa bước vào văn phòng, cảnh tượng đầu tiên đập vào mắt khiến tôi phải dừng bước.

``Korone-chan, em đang làm gì vậy?'' tôi hỏi, nhìn nàng Khuyển đang chạy vòng vòng quanh Aqua như một con chó đuổi đuôi.

``Em đang chơi với Aqua-chan nè\~{}!'' Korone vừa chạy vừa cười toe toét, ``Rảnh quá nên muốn tìm ai đó chơi cùng đó|~{}''

Aqua đứng giữa vòng xoay của Korone, mặt ngơ ngác không hiểu chuyện gì đang xảy ra, mồ hôi túa ra như thác: ``Ơ-Ơ\dots\ Korone-chan? C-Chị làm gì vậy?''

Bỗng nhiên, Korone dừng lại. Em ấy thò tay vào túi và lấy ra một đồng xu sáng loáng (mà sau này tôi mới biết là em ấy mượn từ góc tủ đằng kia).

``Aqua-chan này\~{}'' em ấy bắt đầu lắc lư đồng xu trước mặt nàng Hầu gái, ``Nhìn vào đồng xu này đi\~{}''

``Ể? Đ-Đồng xu?'' đồng nghiệp của em ấy ngây thơ nhìn theo chuyển động của đồng xu, ``Korone-chai định làm gì vậ-''

Tôi chưa kịp ngăn cản thì Akutan đã bắt đầu lảo đảo. Mắt em ấy trở nên mơ màng, người nghiêng ngả như say rượu.

``K-Korone-chan! Dừng lại đi!'' tôi hét lên.

Nhưng Korone-chan chỉ cười khúc khích, bỏ đồng xu vào túi rồi nhảy tót đến chỗ để bánh mì. ``Ôi, đói quá! Ăn chút bánh đã\~{}''

*RẦM!*

Và thế là, nàng Hầu gái đổ xuống sàn như một bao gạo, còn Korone-chan thì vẫn ngồi đó, nhồm nhoàm ăn bánh, như thể chuyện vừa rồi chẳng hề tồn tại.

\begin{figure}[H]
  \centering
  \begin{subfigure}[t]{0.45\textwidth}
    \centering
    \includegraphics[width=8cm]{../../../../Image/Book?/9-5e1.png}
  \end{subfigure}
  \quad
  \begin{subfigure}[t]{0.45\textwidth}
    \centering
    \includegraphics[width=8cm]{../../../../Image/Book?/9-5e2.png}
  \end{subfigure}
  \quad
  \begin{subfigure}[t]{0.45\textwidth}
    \centering
    \includegraphics[width=8cm]{../../../../Image/Book?/9-5e3.png}
  \end{subfigure}
  \caption{Những gì đã diễn ra từ trước đó.}
\end{figure}

\begin{center}
  \uline{Kết thúc hồi tưởng.}
\end{center}

Tôi vừa dứt lời, Kanata-chan đã trợn mắt, quay ngoắt về phía Koro-chan vẫn đang giả vờ vô tội ở đằng kia: ``Tại sao chị lại làm thế chứ!?''

Nàng Khuyển nhún vai, đưa tay chỉ về phía kệ gần đó, nơi mà em ấy vừa lấy đồng xu: ``Thì chị thấy đằng kia có con lắc chứ sao! Thấy cũng vui mắt nên là đem ra thử luôn\~{}''

Tôi nhíu mày, không thể tin được lời đáp đó lại có thể thốt ra một cách hồn nhiên đến vậy.
``Con lắc thì liên quan quái gì đến việc để Aqua-chan bị gục!?''

``Bộ chị chơi lắc siêu xu à!?'' nàng Thiên sứ không chịu được mà bật lại.

Trong lúc tôi còn đang định dọn dẹp lại hiện trường thì em Thiên sứ bất ngờ bùng nổ. Em giận dữ chỉ tay vào nàng Khuyển, hét lên như sấm: ``Tất cả là tại chị đấy! Aqua-chan mà có mệnh hệ gì là lỗi của chị hết! Phải làm sao bây giờ hả!?''

Cơn quát mắng đậm chất ``giáo huấn'' này khiến Koro-chan có phần chột dạ, nhưng vẫn cười cười giả lơ như không nghe thấy. ``Chạy lòng vòng, chơi thôi miên rồi bỏ mặc người ta nằm đó là sao!? Lỡ người ta tắt thở thật thì làm sao hả?''

``Thôi, thôi, Kana-chan ơi, bình tĩnh, đừng nóng\dots'' tôi đang định xen vào giảng hòa hai người thì\dots

*RẮC* *RĂC* *RẮC*

Âm thanh khô khốc vang lên từ trần nhà.

Tôi ngẩng đầu lên, và suýt chút nữa buột miệng chửi thề. ``\vie{Ôi định công mệnh\dots}''

Mảng trần mà Kana-chan vừa sửa chưa được năm phút trước giờ đang rung lắc như muốn từ biệt cuộc đời. Và trước khi tôi kịp hét lên báo động, nó đã rơi xuống với một tiếng *RẦM!* thẳng lên người Aqua vẫn còn đang nằm bất động phía dưới.

\img{9-5f}

Cả ba chúng tôi đồng loạt giật bắn người, mặt tái mét như sáp nến vừa tắt.

``C-Chết toi rồi\dots'' cả Koro-chan và Kana-chan hốt hoảng khi nhìn đồng nghiệp bị những mảng trần lởm chởm kia đè mạnh.

Tôi thì chỉ có thẻ thở dài bất lực, tặc lưỡi: ``Đã bảo là cái trần đó lửng lẳng rồi mà không nghe\dots\ Chậc\dots''

Nhưng đó vẫn chưa phải phần kinh khủng nhất.

Một luồng sáng mờ mờ\dots\ trôi ra khỏi người Aqua-chan. Như thể linh hồn của em ấy vừa\dots\ xuất ra khỏi xác. Một đám mây trong suốt hình hồn ma bé con đang lơ lửng như có định hướng riêng.

\img{9-5g}

Tôi đứng hình, không thốt nổi một câu nào. Koro-chan thì há hốc miệng: ``Cái đó\dots\ là hồn bà ấy đấy à!?'' còn Kana-chan thì rùng mình một cái như vừa thấy ma thật: ``Khoan\dots\ có phải em vừa thấy cái gì bay ra không\dots?''

Tôi gật đầu. Hồn Aqua hít hít không khí rồi nói vọng: ``Mình ngửi thấy mùi sushi ngoài kia\~{}\dots''

Rồi nó lơ lửng bay thẳng ra ngoài hành lang.

Chúng tôi cùng nhau hét lên một tiếng: ``\textbf{Mau bắt bà ấy lại!!}'' và lao ra đuổi theo cái hồn đó.

Trận hỗn chiến chính thức bắt đầu.

Nàng Thiên sứ là người đầu tiên đuổi kịp và cố tóm lấy hồn Aqua, nhưng có cái hồn ấy bị chia thành mấy mảnh nhỏ bay tứ tung.

\gif{9-5h}{1}{37}

``\textbf{Á! Lỡ ép mạnh quá rồi!!}''

Tôi chưa kịp phản ứng thì Korone đã phóng lên tiếp theo, nhảy lên cao vồ một cái như bắt bóng chày. Lần này tưởng như thành công, vì em ấy hét lên: ``\textbf{Rồi! Tóm được rồi!}'' với vẻ mặt cực kỳ mãn nguyện.

Kana-chan vừa thở hổn hển vừa khen: ``Nai-xừ!''

Tôi thì vẫn còn đang vung cái lưới bắt ma thủ công mà Sui-chan đưa cho mượn hồi đợt dọn kho đạo cụ. Nhưng khi nhìn lại thứ Koro-chan đang cầm\dots\ tôi chợt nhận ra có gì đó sai sai.

``Này, K-Koro-chan\dots\ Cái đó là\dots''

Kana-chan hét lên, mặt tái mét: ``\textbf{Cái đó là đuôi của chị mà!!}''

\img{9-5i}

Nàng Khuyển nhìn xuống tay mình. Đúng thật. Là cái đuôi chó của chính mình. Em ấy nhăn mặt vì bắt hớ, còn tôi thì chẳng kịp cảnh báo, cái lưới của tôi đã bay tới, trùm lấy cả hai cô nàng.

``Cái quái gì vậy, Kuon-senpai!?''

``Em có phải là chó hoang, gâu gâu!''

``X-Xin lỗi! Anh bắt nhầm!''

Cả ba chúng tôi vất vả gỡ cái lưới, rồi tiếp tục rượt theo cái hồn kia ra tận ngoài văn phòng.

\begin{dialogue}
  \speak{Kanata} Quay lại đây!
  \speak{Kuon} Đứng lại!
  \speak{Korone} *sủa inh ỏi* Gâu! Gâu gâu gâu!!
\end{dialogue}

Ngay lúc đó, chúng tôi thấy Aki-chan đang đi dạo thong dong, tay cầm quyển sách. Hồn Aqua-chan bay ngang qua, giật luôn đôi bím tóc lơ lửng của em ấy rồi tiếp tục bay đi.

``\textbf{Đôi bím tóc của mình!!}'' Aki-chan thốt lên, mặt đỏ hỏn vì xấu hổ (?).

Tôi ló mặt ra từ phía sau, thở dài: ``Lại nữa à\dots?''

``Đứng lại đó\~{}!!'' em ấy hô to, vậy là thêm một người nhập cuộc. Năm người giờ đang chạy rầm rập như đoàn tàu mini.

\gif{9-5j}{1}{104}

Chưa kịp hoàn hồn, cái hồn kia lại lao vào Marine đang cầm hộp bánh hải sản. Em ấy còn chưa kịp hiểu chuyện quái gì thì bóng hồn kia đã giật phăng cái mũ hải tặc trên đầu.

Senchou rên lên, ôm sát người mình như bị ai đó sàm sỡ: ``An\~{} Mũ của mình!''

Korone xen vào, hớn hở như chơi lễ hội: ``Lẹ lên lẹ lên\~{} gâu gâu\~{}!''

Tôi gục đầu xuống, thở ra vì tiếng ``an\~{}'' không liên quan kia: ``Oắt-đờ-heo đang diễn ra vậy\dots?''

``Trả lại mau!!'' Thêm một người nữa. Sáu người rượt đuổi, mỗi người một tiếng hét, chạy như hỗn chiến đường phố.

\gif{9-5j}{105}{187}

Cuối cùng, khi tới gần Pekora, cái hồn định với tay giật củ cà rốt gắn trên tóc của em ấy\dots\ nhưng không hiểu thế nào lại bị mắc vào đó.

Nàng Thỏ ngơ ngác nhìn cái hồn kia, giật mình: ``Ể? Cái gì thế, peko!?''

\img{9-5k}

Cả nhóm chúng tôi cùng dừng lại, ánh mắt cùng đổ dồn vào cái hồn đang kẹt trong lọn tóc ấy. Tôi lập tức chỉ tay chớp thời cơ: ``Đứng yên ở đó, Peko-chan! \textbf{Bắt lấy}!!''

Ngay lập tức, năm người kia lao vào như bầy thú hoang, tay chân loạn xạ, túm lấy cái hồn Aqua khi nó còn đang lơ lửng không hiểu chuyện gì xảy ra.

Tôi quăng lưới một lần nữa. lần này là chuẩn xác\dots

``Âu sịt\dots''

\dots chuẩn xác tới mức bắt trọn cả năm người con gái luôn.

*PHỊCH!*

Hồn Aqua-chan, cùng với Kana-chan, Koro-chan, Aki-chan, Marine-chan và cả Peko-chan nằm gọn trong cái lưới khổng lồ ấy như một tổ sushi cuộn.

Một hồi sau, sau khi lồm cồm bò ra khỏi đống lưới, trên tay họ vẫn giữ được cái hồn đã vật vờ kia.

Tôi chưa kịp thở thì ba người trong số họ đồng loạt mắng tôi tơi bời trong khi còn đang tìm cách thoát khỏi cái lưới bắt ma kia.

\begin{dialogue}
  \speak{Kanata} Anh làm cái trò gì vậy hả!?
  \speak{Marine} Tự nhiên thả lưới lên đầu phụ nữ là sao hảaa\~{}?
  \speak{Pekora} Peko\dots\ Đau đầu quá trời luôn peko!!
\end{dialogue}

Tôi chắp tay xin lỗi rối rít, cúi đầu liên tục. May mà Aki-chan và Koro-chan thì chỉ cười xòa cho qua, thậm chí còn đập tay nhau:
\begin{dialogue}
  \speak{Aki} Lâu rồi mới vận động đã vậy đó\~{}
  \speak{Korone} Hihi, vui dễ sợ luôn á\~{}
\end{dialogue}

Tôi thở phào nhẹ nhõm. Nhưng thề luôn, từ giờ tôi sẽ không bao giờ để cho Kana-chan sửa trần nữa.
%==========================================TRUYỆN PHỤ=========================================

\omk

Chúng tôi đem cái hồn của Akutan - lúc này đang loi nhoi trên tay của nàng Thuyền trưởng - về lại văn phòng như thể vừa bắt được một sinh vật huyền thoại. Nó bay vòng vòng rồi đâm đầu vào tóc Marine, sau đó tự mình xoắn thành hình xoắn ốc như thể đang giả làm sushi cá ngừ. Marine-chan phải giữ chặt nó lại bằng cả hai tay: ``Nào, đứng yên đi! Bộ chị muốn làm topping luôn hả!?''

Tôi nhìn nó, đầu hơi nghiêng nghiêng. Dù nó trông như một cục bông phát sáng lập lòe, có phần lộn xộn\dots\ nhưng vẫn toát ra nét gì đó đáng yêu.

``Nếu nhìn kỹ thì\dots\ Aqua-chan dạng hồn cũng dễ thương đấy chứ,'' tôi nói, tay chống cằm suy tư như một nhà nghiên cứu khoa học kỳ quặc. ``Kiểu trông như mascot ấy.''

\img{9-5l}

Những người còn lại nhìn nhau. Bầu không khí ngập ngừng trong một giây, rồi mọi người cùng nhau gật đầu đồng tình, thậm chí có người còn bị ấn tượng bởi cái hồn này.

``Nghe cũng hợp lý phết,'' Aki0chan mỉm cười, tay chống hông. ``Trông như một cục bông đang lơ lửng trên trời vậy.''

``Ý của chị là cục mây à, peko.''

``Sao ta không lấy nó làm linh vật chính thức cho \textit{hologra} nhỉ?'' Marin-chan búng tay cái tách, mắt sáng rỡ như tìm được kho báu.

Nàng Dị giới reo lên: ``Ý hay đấy!''

``Ể!? Có ổn không vậy!?'' Kana-chan cau mày, nửa lo lắng, nửa nghi ngờ, trong khi cái hồn kia vẫn đang múa may quay cuồng trên tay Marine như đang luyện ba-lê trên không.

``Tôi thấy nó lấn cấn thế nào ấy, peko\dots'' Pekora gãi đầu, tai thỏ giật nhẹ.

``Anh thấy cũng được mà,'' tôi nhún vai.

Nàng Thỏ và nàng Thiên sứ liền nhìn sang tôi, nghiêng đầu hỏi: ``Nhưng bằng cách nào?''

Tôi rút máy ảnh trong túi áo ra, lắc nhẹ. ``Chụp một tấm ảnh cái hồn này, sau đó dựng lại bằng model 3D. Dễ òm.''

Không đợi thêm giây nào, tôi bấm máy một cái tách. Ngạc nhiên thay, nó hiện lên nguyên vẹn trong ảnh, đủ chi tiết để đem đi dựng. Mọi người cùng xúm vào xem.

Aki-chan thốt lên: ``Trông dễ thương ghê!''

Marine thì cười khúc khích, đưa ngón tay dí vào màn hình: ``Nhỏ nhỏ, lơ lửng\~{} Nhìn cứ như kẹo bông gòn biết bay vậy nha\~{}''

Pekora giật giật ống tay áo tôi: ``Làm sao mà anh chụp được vậy, peko? Linh hồn thì đáng lẽ đâu có hiện hình được mà, đúng không!?''

Tôi nháy mắt: ``Đây là \hll\ mà. Cái quái gì cũng xảy ra được.''

\ct

Sau một hồi cãi nhau, bàn bạc, rồi thử nghiệm vài ba lần gần như thành công (mà thực chất là hồn Aqua lại bay ra vì hắt hơi đúng lúc), cuối cùng chúng tôi cũng nhập được cái hồn kia trở lại vào người em ấy. Tôi dùng một thiết bị chuyên dụng dạng ``máy hút chân không linh hồn'' (do tôi và Miharu chế), còn Kana-chan thì vừa giữ đầu, vừa thì thầm mấy câu cầu nguyện nghe như đang làm lễ trừ tà.

Một luồng sáng mờ mờ lóe lên, rồi thân thể Aqua giật nhẹ. Mọi người nín thở.

Cô nàng khẽ rên một tiếng, rồi cựa quậy. Trên đầu cô lúc này vẫn còn dính nguyên mũ thuyền trưởng của Marine-chan, đôi bím tóc vàng óng ánh của Aki-chan, và củ cà rốt lủng lẳng từ lọn tóc của Peko-chan.

Một hình dạng\dots\ trông khá khó diễn tả, như thể một cô nàng không biết phải hóa trang vào ai, để rồi ``chơi lớn'', lấy một phần nhân vật một chút và tự biến mình thành một sản phẩm lỗi.

``Ể!? Ủa\dots? Mình đang ở đâu\dots?'' Aqua thều thào, mắt đảo một vòng như thể vừa tỉnh dậy giữa giấc mơ sau một đêm dài.

\begin{figure}[H]
  \centering
  \begin{subfigure}[t]{0.45\textwidth}
    \centering
    \includegraphics[width=8cm]{../../../../Image/Book?/9-5m1.png}
    \caption{Aqua tinh lại, với những phụ kiện vẫn còn trên đầu.}
  \end{subfigure}
  \quad
  \begin{subfigure}[t]{0.45\textwidth}
    \centering
    \includegraphics[width=8cm]{../../../../Image/Book?/9-5m2.png}
    \caption{Năm cô gái đứng nhìn tình hình của cô với từng biểu cảm khác nhau.}
  \end{subfigure}
\end{figure}

Khi em ấy mở mắt hẳn và ngồi dậy, thứ đầu tiên đập vào mắt mình là sáu người đang đứng thành hàng trước mặt mình, trong đó có ba gương mặt của ``chủ nhân hợp pháp'' các món phụ kiện giờ đang trên đầu của chính em ấy.

Marine-chan khoanh tay, nhướng một bên mày, chống hông.

Peko-chan thì khoanh tay thở dài: ``Cái củ đó để ăn, không phải để đội đâu, peko.''

Aki-chan, tất nhiên, vẫn cười nhẹ: ``Trông cũng hợp với cậu đó nha.''

Tôi thì đang tủm tỉm, tay chìa ra hộp sushi cao cấp đã đặt từ sáng: ``Tỉnh rồi hả, Akutan? Làm chút sushi không?''

Nàng Hầu gái vẫn còn ngơ ngác, miệng há hốc nhìn tôi như thể tôi là người ngoài hành tinh, nhưng rồi, bản năng mạnh hơn lý trí, em ấy cầm lấy hộp sushi và mở nắp như bị điều khiển. Tay còn lại đưa lên đầu gỡ từng món ra, vừa lẩm bẩm vừa lần lượt trả lại cho từng người, những người hơi xấu hổ nhưng vẫn cố giữ thể diện.

Tôi tranh thủ bước lại gần và nhỏ giọng: ``Lần sau nhớ tránh xa mấy cái lỗ trên trần nhà nhé.'' Hiển nhiên, đáp lại tôi chỉ là một dấu hỏi to đùng lửng lơ trên đầu, không hiểu tôi đang nói đến cái gì.

Ngay sau đó, hai cô gái còn lại bất ngờ cúi rạp người xuống: ``Xin lỗi vì đã gây rắc rối!!!''

Kana-chan là người cúi đầu trước, tay còn giữ chặt đầu Koro-chan ép xuống, vì rõ ràng nàng Khuyển kia vẫn đang toe toét như thể chẳng có gì xảy ra. ``A, đau đau đau- nhưng mà em có làm gì sai đâu mà!?'' Koro-chan vẫn cố nhặng xị lên dù mặt dính sát sàn.

\img{9-5n}

Bốn cô gái còn lại nhìn nhau như thể đang chứng kiến một nghi lễ lạ đời. ``Họ xin lỗi cái gì vậy, (peko)\dots?''

Tôi khẽ thở dài, đứng giữa nhóm và bắt đầu giải thích lại toàn bộ chuỗi sự kiện điên rồ vừa rồi: từ mảng trần rơi, đến vụ linh hồn Akutan bay đi săn sushi, đến màn đuổi bắt khiến cả dãy văn phòng suýt nữa thành chiến trường.

Sau khi nghe xong, bốn người còn lại chỉ biết trố mắt và rồi cũng hiểu dần đầu đuôi câu chuyện.

``\dots Thế nê mới ra nông nỗi này,'' tôi kết luận.

Nàng Thuyền trưởng khoanh tay: ``Tôi nên biết là có liên quan tới Kanata và Korone mà.''

Nàng Hầu gái vẫn đang ăn món sushi, ngước mắt ngơ ngác nhìn mọi chuyện: ``Ể? Vậy là chị vừa mới chết lâm sàng à?''

``Có thể hiểu như vậy,'' tôi nhún vai.

Tôi vỗ vai nàng Thiên sứ, nhắc nhở thêm một câu: ``À mà này, Kanata-chan\dots\ anh nghĩ em cũng nên xin lỗi cả A-san nữa.''

Nghe tới đây, cô nàng Thế hệ thứ Tư như bị sét đánh. Em ấy sực nhớ lại cái trần nhà ban nãy được sửa một cách tạm bợ và nửa vời, và giờ chắc phải làm lại toàn bộ.

Em ấy chán nản thở dài: ``Biết rồi\dots\ Em sẽ tự bỏ tiền túi ra sửa.''

Tôi gật đầu: ``Tốt. Anh sẽ báo lại cho A-san. Giờ thì đứa nào dìu chị ấy xuống phòng y tế nhỉ?''

Cả nàng Dị giới và nàng Thuyền trưởng giơ tay xung phong. ``Được rồi. Mấy đứa đi đi.''

Ở phía sau, Koro-chan ngước đầu dậy, mắt sáng lên: ``Vậy là khỏi bị trừ lương hả!?''

``Không em.''

``Ê, khoan khoan, sao em cũng bị trừ vậy!?''

``Vì em không giúp được gì,'' tôi đáp gọn lỏn, và cả phòng bật cười. Nàng Khuyển không thể làm gì hơn ngoài việc bĩu môi.

\ct

Bầu không khí sau cơn hỗn loạn cuối cùng cũng dịu lại. Aqua-chan đã hoàn hồn\dots\ theo cả nghĩa đen lẫn nghĩa bóng, và giờ đang ngồi nhai miếng sushi cuối cùng.

Kana-chan thì vừa dùng dụng cụ hàn chuyên dụng để vá lại phần trần nhà mà em ấy đã lỡ tay phá hỏng, miệng vừa than: ``Trời ơi\dots\ tại sao cái trần này lại cứng thế\dots\ Vai em mỏi quá\dots''

Tôi đứng ở dưới, khoanh tay đáp lại: ``Đừng có than nữa. Anh tưởng cơ tay của em khỏe lắm mà?''

``Lực nắm của em mới mạnh, chứ không phải cơ tay của em!''

``Rồi, rồi, biết thế. Làm nốt đi, làm cho cẩn thận vào, chứ đừng có cẩu thả như vừa nãy đâu đấy.''

Đến khi cô lén định hàn đại cho xong, tôi chỉ nhích nhẹ khóe miệng, và cô lập tức siết chặt tay cầm như thể vừa được tiêm doping.

``Chết rồi, làm xấu là bị A-chan mắng thật luôn á\dots'' em ấy méo mặt.

Cuối cùng, nhờ vào bộ công cụ mà tôi mang theo, phần trần cũng được khôi phục nguyên trạng. Tôi gật gù, xem như xong xuôi.

Lúc ấy, tôi chợt nhớ đến một việc quan trọng. Quay sang hai người còn lại đang không đụng tay đụng chân, tôi hỏi: ``À mà\dots\ trước khi hai em rời khỏi nhà, đã dọn lại cho sạch sẽ chưa?''

``Rồi ạ!'' họ đồng thanh đáp.

``Em có gom mấy cái gối thành một đống đó, được không?'' nàng Khuyển nhanh nhảu giơ tay lên.

``Miễn là làm thế nào để tầng ba sạch sẽ là được,'' tôi đáp.

``Em với Mio-chan cũng đã quét dọn với lau qua tầng ba nhà anh rồi đó,'' đến lượt nàng Hầu gái báo cáo.

``Vậy là ổn rồi. Cảm ơn mấy đứa.''

Như vậy, sau bao nhiêu rối ren, náo loạn và ``sự cố hồn lìa khỏi xác,'' thật sự thì\dots\ cuộc sống làm việc ở trong công ty này đã quay trở về bình thường.

Vẫn là những trò con bò ấy, vẫn là những sự cố làm cho chúng tôi thót tim đấy, vẫn như thường ngày, không có gì thay đổi.

Nhưng này, như thế mới chính là \hll\ chứ, đúng không?
\end{document}